\clearpage
\VerdierPhysicalPageAnchor{0137}%
\begingroup
\small
\setlength{\emergencystretch}{1.5em}%
\setlength{\abovedisplayskip}{1pt plus 1pt minus 1pt}%
\setlength{\belowdisplayskip}{1pt plus 1pt minus 1pt}%

\VerdierSemanticAnchor{ch2-2-2-2-sm2-completion-diagram}
\noindent en un diagramme commutatif :

\begin{center}
  \begin{tikzpicture}[>=stealth,line width=.38pt,x=1.30cm,y=1.45cm,
    scale=.92,transform shape]
    \node (Zpp) at (0,2.20) {$Z''$};
    \node (Xp)  at (-1.25,1.10) {$X'$};
    \node (Yp)  at ( 1.25,1.10) {$Y'$};
    \node (Y)   at (0,0) {$Y$};
    \draw[->] (Zpp) -- node[above left=-1pt] {$t'\in S$} (Xp);
    \draw[->] (Zpp) -- node[above right=-1pt] {$m'$} (Yp);
    \draw[->] (Xp)  -- node[below left=-1pt] {$m$} (Y);
    \draw[->] (Yp)  -- node[below right=-1pt] {$t\in S$} (Y);
    \node at (2.50,-.05) {$.$};
  \end{tikzpicture}
\end{center}

\VerdierSemanticAnchor{ch2-2-2-2-composite-roof-d3}
\noindent On obtient alors un diagramme
$D_3=X\xleftarrow{st'}Z''\xrightarrow{pm'}Z$, où le composé $st'$ appartient
à $S$ (SM1).
\VerdierSemanticAnchor{ch2-2-2-2-composite-class-independence}%
On vérifie que la classe $\overline D_3$ de $D_3$ dans
$\operatorname{Hom}_{\mathcal C(S^{-1})}(X,Z)$ ne dépend que des classes
$\overline D_1$ et $\overline D_2$.
\VerdierSemanticAnchor{ch2-2-2-2-composition-conclusion}%
Le morphisme $\overline D_3$ est le
composé dans $\mathcal C(S^{-1})$ des morphismes $\overline D_1$ et
$\overline D_2$.

\VerdierSemanticAnchor{ch2-2-2-3}
\smallskip
\noindent\textbf{Remarque 2.2.3.}
\VerdierSemanticAnchor{ch2-2-2-3-a}%
a) Il résulte de la quatrième assertion du théorème (2.2.2) que la catégorie
$\mathcal C(S^{-1})$ munie du foncteur
$Q\colon\mathcal C\longrightarrow\mathcal C(S^{-1})$ est solution d'un
problème universel :
\VerdierSemanticAnchor{ch2-2-2-3-a-universal-problem}%
le foncteur $Q$ transforme les morphismes de $S$ en isomorphismes, et tout
foncteur de $\mathcal C$ dans une catégorie $\mathcal A$ qui transforme les
morphismes de $S$ en isomorphismes se factorise d'une manière unique par $Q$.
\VerdierSemanticAnchor{ch2-2-2-3-a-unique-isomorphism}%
La catégorie $\mathcal C(S^{-1})$ munie du foncteur $Q$ est donc déterminée
à isomorphisme unique près par cette propriété.

\VerdierSemanticAnchor{ch2-2-2-3-b}
\par\noindent b) Soient $\mathcal C$ une catégorie et $S$ un ensemble de
morphismes de $\mathcal C$ qui ne soit pas nécessairement un système
multiplicatif de $\mathcal C$.
\VerdierSemanticAnchor{ch2-2-2-3-b-general-localization}%
Il existe toujours une catégorie $\mathcal C(S^{-1})$ et un foncteur
$Q\colon\mathcal C\longrightarrow\mathcal C(S^{-1})$ qui transforme les
morphismes de $S$ en isomorphismes, tel que tout foncteur de $\mathcal C$
dans une catégorie $\mathcal A$ transformant les morphismes de $S$ en
isomorphismes se factorise d'une manière unique par $Q$.
\VerdierSemanticAnchor{ch2-2-2-3-b-multiplicative-comparison}%
Lorsque $S$ est un système multiplicatif de $\mathcal C$, le théorème
(2.2.2) donne une description simple de la catégorie $\mathcal C(S^{-1})$ ;
il donne de plus des propriétés supplémentaires du foncteur $Q$ (assertions
b) et c)) qui ne sont pas vraies en général.

\VerdierSemanticAnchor{ch2-2-2-4}
\smallskip
\noindent\textbf{Corollaire 2.2.4.}\quad
\begingroup\itshape
Soient $\mathcal C$ une catégorie et $S$ un système multiplicatif de
$\mathcal C$.

\VerdierSemanticAnchor{ch2-2-2-4-a}
\par\noindent a) L'ensemble $S$ est un système multiplicatif de
$\mathcal C^\circ$, la catégorie opposée à $\mathcal C$.
\VerdierSemanticAnchor{ch2-2-2-4-a-opposite-system}%
On a un isomorphisme canonique de catégories :

\VerdierSemanticAnchor{ch2-2-2-4-a-canonical-isomorphism}
\[
  \bigl(\mathcal C^\circ(S^{-1})\bigr)^\circ
  \longrightarrow \mathcal C(S^{-1})\,.
\]

\VerdierSemanticAnchor{ch2-2-2-4-a-hom-colimit-isomorphism}
\noindent Cet isomorphisme induit en particulier des isomorphismes :
\[
  \operatorname{Hom}_{\mathcal C(S^{-1})}(X,Y)
  \longrightarrow
  \varinjlim_{Y\backslash S}\operatorname{Hom}_{\mathcal C}(X,\mathord\cdot)\,.
\]
\endgroup

\VerdierSemanticAnchor{ch2-2-2-4-p138-boundary}
\endgroup
